\section{What the occupation law determines}\label{sec:reconstruction}

The forward theory says which face carries the mode for a given chain. This
section asks which coarse functionals of the law of the occupation counts
determine the generator, and how much they lose. We compare three of them, from
fine to coarse. The first is the law of the counts itself, at every $N$ and $T$.
The second is its logarithmic scale on single compositions, which by the
stratified local law (Theorem~\ref{thm:stratified}) is the Donsker and Varadhan
functional $J_V$. The third is a finite label, the modal support of the chain
killed at known rates (Section~\ref{sec:observation}). The organizing result is
that the second and the third carry the same information, namely the principal
minors of $-Q$ (Corollary~\ref{cor:dv-fiber}). Theorem~\ref{thm:nonreversible}
describes this information. It is $Q$ up to time reversal when $Q$ has no cut,
and a finite fiber of generators in general. For a reversible chain on a known
graph it is all of $Q$, and finitely many boundary coordinates of a design fixed
in advance recover it (Section~\ref{sec:reversible-reconstruction}).

\subsection{The killed chain and its modal support}\label{sec:observation}

Fix a field $v\in\R^d$ and a time scale $\tau>0$, put $T=\tau L$ and $h=T/N$,
and weight each composition by $e^{Tv\cdot n/N}$. Along a path the weight is
$\prod_te^{hv_{X_t}}$. If $v\le0$, this is the probability that the chain
survives when it is killed at each visit to $i$ with probability
$1-e^{hv_i}$. So $P_N(n)e^{Tv\cdot n/N}$ is the probability of surviving with
composition $n$, and its maximizers are the modes of the composition of the
surviving chain, since a normalization does not move a mode. On a face the extra
killing shifts the diagonal of the killed generator, $Q_F+\operatorname{diag}v_F$
in place of $Q_F$. For a face $F$ let
\[
 \lambda_F(v)=-\varrho\bigl(Q_F+\operatorname{diag}v_F\bigr),\qquad
 \Phi_{\tau,v}(F)=|F|-1+\tau\lambda_F(v),
\]
so $\lambda_F(0)=\lambda_F$ and $\Phi_{\tau,0}=\Phi_\tau$. If $F$ is reducible,
$Q_F$ is block triangular in the order of its strongly connected components, and
$\lambda_F(v)$ is the least of the numbers $\lambda_C(v)$ over these components $C$.

The shift $v\mapsto v+c\ones$ multiplies the weighted law by the constant
$e^{Tc}$, and since $\lambda_F(v+c\ones)=\lambda_F(v)-c$ it moves every score
$\Phi_{\tau,v}(F)$ by the same $-\tau c$. So it leaves every winner unchanged,
and every field may be taken nonpositive. A nonpositive field is a known killing
rate, and $\lambda_F(v)$ is the exit rate of $F$ for the killed chain. The
observation is then the support of the most likely composition of the surviving
trajectory, for a chain with known extra killing watched up to the horizon
$T=\tau L$. Throughout this section the field and the time scale are
calibrated, that is, $v$ and $\tau$ are known in the units of $Q$.

The survival probability itself, $\sum_nP_N(n)e^{Tv\cdot n/N}$, is a finer
observation. For irreducible $Q$ it equals $e^{-T\lambda_V(v)}$ up to a bounded
factor, so it gives $\lambda_V$ on all fields, and by
Lemma~\ref{lem:perron-branch} below it already determines every principal minor
of $-Q$, with no descent through the faces. The modal-support label is coarser.
It does not require measuring a probability, only the support of the most likely
composition, and it is the datum that the forward theory describes. We show that
it loses nothing.

The controlled score law below says that the face that maximizes the weighted
law has the least score $\Phi_{\tau,v}$, up to $O(\log L)$ in the exponent. A
face $F$ wins strictly at $(\tau,v)$ if it is the unique minimizer of
$\Phi_{\tau,v}$ over the admissible faces. The modal-support map is
$(\tau,v)\mapsto F$ on the strict winning regions, and the data are the
boundaries between adjacent strict regions, the controlled boundaries. We do not assume that a face tied
in $\Phi$ is observed as a maximizer on the boundary. For two states a factor
$\sqrt{L}$ decides a tie, and we do not use it.

\begin{proposition}[the controlled score law]\label{prop:controlled-score}
Fix $Q$, a start $\mu$ with full support, a field $v$ and a compact interval of
$\tau$ in $(0,\infty)$. For every admissible face $F$,
\[
 \log\max_{\supp n=F}\bigl\{P_N(n)e^{Tv\cdot n/N}\bigr\}
 \le-\Phi_{\tau,v}(F)\,L+O(\log L),
\]
with equality up to $O(\log L)$ if $F$ is irreducible, and for a singleton the
maximum is exactly $\mu_i(1-hq_i)^{N-1}e^{Tv_i}$. A reducible admissible face
$F$ has $\Phi_{\tau,v}(F)\ge\Phi_{\tau,v}(C)+1$ for one of its strongly
connected components $C$. Hence if $F$ is the unique minimizer of
$\Phi_{\tau,v}$ over the admissible faces, then $F$ is irreducible and the
support of every maximizer of the weighted law is $F$ for all large $N$.
\end{proposition}

\begin{proof}
The proposition is the field form of Theorem~\ref{thm:first}(ii) and (iii). Let
$k=|F|$. For large $N$ the whole interval of $\tau$ gives $T_0\le T=\tau L\le N^{1/3}$,
and $\log T=O(\log L)$. The only composition with support $\{i\}$ is $Ne_i$, so
by~\eqref{eq:vertexmax} the maximum is exactly $\mu_i(1-hq_i)^{N-1}e^{Tv_i}$, whose
logarithm is $-T(q_i-v_i)+O(1)=-\Phi_{\tau,v}(\{i\})L+O(1)$. For $k\ge2$,
Remark~\ref{rem:first-field} gives the upper bound of Theorem~\ref{thm:first}(ii)
for $M_F(T,v)$ with $\lambda_F(v)$ in place of $\lambda_F$. It uses~\eqref{eq:upper-field}
on each block of a path, so it holds for reducible faces too, and it reads
$\log M_F(T,v)\le-(k-1)L-T\lambda_F(v)+(k-1)\log T+\kappa=-\Phi_{\tau,v}(F)L+O(\log L)$.
If $F$ is irreducible, it satisfies (H) because $\mu$ has full support, and~\eqref{eq:logmax}
of Theorem~\ref{thm:facemax} with $\mathcal V=\{v\}$ gives the matching lower bound.

Let $F$ be reducible and admissible. Then $\lambda_F(v)$ is the least
$\lambda_C(v)$ over the components $C$ of $F$, and a component $C$ where it is attained has
$|C|\le k-1$. It is admissible by Lemma~\ref{lem:admissible}(b), and
$\Phi_{\tau,v}(F)=\Phi_{\tau,v}(C)+k-|C|\ge\Phi_{\tau,v}(C)+1$. So a unique minimizer $F$
of $\Phi_{\tau,v}$ is irreducible. Let $\varepsilon>0$ be the least gap
$\Phi_{\tau,v}(G)-\Phi_{\tau,v}(F)$ over the other admissible faces $G$. Their weighted
maxima are at most $N^{-\Phi_{\tau,v}(F)-\varepsilon+o(1)}$, that of $F$ is at least
$N^{-\Phi_{\tau,v}(F)-o(1)}$, and a face that is not admissible carries no mass. There are
finitely many faces, so for large $N$ every maximizer has support $F$.
\end{proof}

The upper bound holds up to the boundary of the face, and this is what lets a
fixed field act on all compositions at once.

\begin{definition}[controlled boundaries]\label{def:controlled}
Two generators have the same controlled boundaries if, for every start of full
support and every field used below, they have the same strict winner at every
point $(\tau,v)$ outside a closed nowhere dense set. The fields used are the
pair fields of Section~\ref{sec:nonreversible-reconstruction}, with a field
$\delta$ on one state, $0$ on another and $-\bar v$ elsewhere, and the fields
$(v_F,-\bar v)$ of Lemma~\ref{lem:exposure}, equal to $v_F$ on a face $F$ and to
$-\bar v$ off it, with $\bar v$ large.
\end{definition}

The ties form a closed nowhere dense set. A minimizer of $\Phi_{\tau,v}$ is irreducible, and for
irreducible faces $F\ne H$ the difference $\Phi_{\tau,v}(F)-\Phi_{\tau,v}(H)$
vanishes for at most one $\tau$ at fixed $v$ if $|F|\ne|H|$. If $|F|=|H|$ and
$j\in F\setminus H$, it vanishes for at most one $v_j$ when the other
coordinates are fixed, since $\lambda_H(v)$ does not depend on $v_j$ while
$\partial\lambda_F/\partial v_j=-(l\circ r)_j<0$ for the Perron vectors of
$Q_F+\operatorname{diag}v_F$ with $l\cdot r=1$. So the set where the least score
is attained twice is closed and nowhere dense. Off it
Proposition~\ref{prop:controlled-score} identifies the observed support with the
strict winner for all large $N$, and the ties themselves are never used.

The organizing result of the section compares the second and the third
functional of the introduction.

\begin{corollary}[the Donsker and Varadhan functional and the modal supports]\label{cor:dv-fiber}
For two irreducible generators $Q$ and $Q'$ the following are equivalent.
\begin{enumerate}[(i)]
\item They have the same Donsker and Varadhan functional $J_V$.
\item They have the same controlled boundaries.
\item They have the same principal minors, that is, they lie in the same fiber
of Theorem~\ref{thm:nonreversible}(b).
\end{enumerate}
\end{corollary}

The proof is given after Theorem~\ref{thm:nonreversible}. The equivalence of
(i) and (iii) uses only the variational formula, Lemma~\ref{lem:RF-irreducible}
and Lemma~\ref{lem:perron-branch}, and (ii) is the strengthening to a single
label. The law of the counts itself does better for a generic start. From the
stationary start the law of the composition is the same under $Q$ and under its
time reversal $Q^*$, for every $N$ and $T$ with $Tq_{\max}\le N$
(Theorem~\ref{thm:nonreversible}(f)), so from that start no count statistic
separates them. A start outside a finite union of proper subspaces separates
every member of the fiber through the amplitudes $K_V(v)$.

\subsection{The reversible case}\label{sec:reversible-reconstruction}

For a reversible chain the information is all of $Q$, and an explicit finite
experiment reads it. Let $Q$ be reversible and irreducible with stationary law
$\pi$, on a known connected graph $G=(V,E)$ with $m=|E|$ edges, where
$\{i,j\}\in E$ exactly when $q_{ij}>0$. Put $w_{ij}=q_{ij}q_{ji}$ for the edge
products and
\[
 S=D_{\sqrt\pi}(-Q)D_{\sqrt\pi}^{-1},\qquad S_{ii}=q_i,\qquad S_{ij}=-\sqrt{w_{ij}}\ \ (\{i,j\}\in E),
\]
the symmetrized killed matrix. We assume known constants $R,b>0$ with
$\|S\|_2\le R$ and $\sqrt{w_{ij}}\ge b$ on every edge. The field and the time
are calibrated, that is, $v$ and $\tau$ are known in the units of $Q$.

Two facts are used throughout. First, $S$ is symmetric and similar to $-Q$,
whose eigenvalues have nonnegative real parts, so $S$ is positive
semidefinite. Hence $0\le q_k\le R$ for every $k$, $\sqrt{w_{kl}}\le R$ on every
edge, and $|q_k-q_l|\le R$. Second, for a face $F$ on which $G$ induces a
connected graph, $-(Q_F+\operatorname{diag}v_F)$ is similar to
$S_F-\operatorname{diag}v_F$ by the restriction of $D_{\sqrt\pi}$, so
\begin{equation}\label{eq:lambda-sym}
 \lambda_F(v)=\lambda_{\min}\bigl(S_F-\operatorname{diag}v_F\bigr),
\end{equation}
as in Proposition~\ref{prop:variational}. Since $Q$ has symmetric support, the
admissible faces are exactly these connected faces
(Lemma~\ref{lem:admissible}(c)).

The design uses, for a retained pair $(i,j)$, the field
\[
 v_i=\delta,\qquad v_j=0,\qquad v_k=-\bar v\ \ (k\ne i,j),
\]
which suppresses every other state, with the constants
\[
 \bar\delta=2R,\qquad \tau_0=\frac1{2R},\qquad U=\frac{8R}{b^2},\qquad \bar v=4R+4UR^2,
\]
and the ranges $|\delta|\le3R$ and $\frac1{4R}\le\tau\le U$.

\begin{lemma}[suppression]\label{lem:suppression}
Under the design for the pair $(i,j)$, every irreducible face that contains a
suppressed state has score at least $\frac34$ larger than the best face inside
$\{i,j\}$, uniformly on the ranges of $\delta$ and $\tau$. Hence the strict
winner is one of $\{i\}$, $\{j\}$ and $\{i,j\}$, and the last only if
$\{i,j\}$ is an edge.
\end{lemma}

\begin{proof}
Let $F$ be connected and let $J=F\setminus\{i,j\}$ be nonempty.

\emph{Wholly suppressed faces.} If $F\cap\{i,j\}=\emptyset$, then
$S_F-\operatorname{diag}v_F=S_F+\bar vI\ge\bar vI$, so $\lambda_F(v)\ge \bar v$
by~\eqref{eq:lambda-sym}. The singleton $\{j\}$ has score $\tau q_j\le\tau R$.
Hence $\Phi_{\tau,v}(F)-\Phi_{\tau,v}(\{j\})\ge\tau(\bar v-R)\ge(3R+4UR^2)/(4R)\ge\frac34$.

\emph{Mixed faces.} Let $G_0=F\cap\{i,j\}$ be nonempty. In the order $G_0$, $J$,
\[
 S_F-\operatorname{diag}v_F=\begin{pmatrix}B&E\\E^\top&S_J+\bar vI\end{pmatrix},\qquad
 B=S_{G_0}-\operatorname{diag}v_{G_0},\quad E=S[G_0,J].
\]
Here $\|E\|_2\le R$ and $S_J\ge0$. Let $e=\lambda_{\min}(B)$, so
$|e|\le\|S_{G_0}\|_2+|\delta|\le4R$. The test vector $(x,0)$, with $x$ a unit
eigenvector of $B$ for $e$, gives $\lambda_F(v)\le e$. For a unit vector $(x,y)$,
\[
 \Bigl\langle\binom xy,(S_F-\operatorname{diag}v_F)\binom xy\Bigr\rangle
 \ge e|x|^2-2R|x||y|+\bar v|y|^2\ge\lambda_{\min}\begin{pmatrix}e&-R\\-R&\bar v\end{pmatrix}
 \ge e-\frac{R^2}{\bar v-e},
\]
by $\sqrt{s^2+R^2}\le s+R^2/(2s)$ with $s=(\bar v-e)/2>0$. Since
$\bar v-e\ge \bar v-4R=4UR^2$, we get $e-1/(4U)\le\lambda_F(v)\le e$. Now
$e=\lambda_{G_*}(v)$ for an irreducible face $G_*\subseteq\{i,j\}$. If $G_0$ is
connected, take $G_*=G_0$ and use~\eqref{eq:lambda-sym}. If $G_0=\{i,j\}$ is not
an edge, $B$ is diagonal and $G_*$ is the singleton with the smaller diagonal
entry. In both cases $|G_*|\le|F|-1$, so
$\Phi_{\tau,v}(F)-\Phi_{\tau,v}(G_*)\ge1-\tau/(4U)\ge\frac34$ as $\tau\le U$.

So the unique minimizer of $\Phi_{\tau,v}$, if there is one, is an irreducible
face inside $\{i,j\}$, and Proposition~\ref{prop:controlled-score} identifies it
with the strict winner.
\end{proof}

\begin{theorem}[reversible reconstruction]\label{thm:reversible-reconstruction}
Let $Q$ be reversible and irreducible on the known connected graph $G$ with
$\|S\|_2\le R$ and $\sqrt{w_{ij}}\ge b$ on edges, and let the start have full
support.
\begin{enumerate}[(a)]
\item (Exit differences.) For each $i\ne1$, under the design for the pair
$(i,1)$ at $\tau=\tau_0$, the strict winner changes from $\{1\}$ to $\{i\}$ at
$\delta=x_i=q_i-q_1$, and $|x_i|\le R$. Put $d_i=x_i$ and $d_1=0$.
\item (Products.) For each edge $\{i,j\}$, under the design for $(i,j)$ at
$\delta=\bar\delta$, the strict winner changes from a singleton to the pair at
\[
 \tau=t_{ij}=\frac{\sqrt{(\bar\delta-\Delta_{ij})^2+4w_{ij}}+\bar\delta-\Delta_{ij}}{2w_{ij}},\qquad
 \Delta_{ij}=d_i-d_j,
\]
with $1/R\le t_{ij}<U/2$, and $w_{ij}=(\bar\delta-\Delta_{ij})/t_{ij}+t_{ij}^{-2}$.
\item (Reconstruction.) Let $B$ be the symmetric matrix with diagonal $d_i$,
entries $-\sqrt{w_{ij}}$ on the edges and zero elsewhere, let
$\lambda=\lambda_{\min}(B)$ and let $\phi$ be its positive unit eigenvector. Then
\[
 S=B-\lambda I,\qquad Q=-D_\phi^{-1}SD_\phi,\qquad \pi_i=\phi_i^2 .
\]
So the $m+d-1$ limiting boundary coordinates $(d_i)_{i\ne1}$ and
$(t_{ij})_{\{i,j\}\in E}$ determine $Q$, and the design is fixed by $R$, $b$ and
the labels before any boundary is observed.
\end{enumerate}
\end{theorem}

\begin{proof}
By Lemma~\ref{lem:suppression} and Proposition~\ref{prop:controlled-score}, at
each point of the ranges the strict winner is the unique minimizer of
$\Phi_{\tau,v}$ among the irreducible faces inside the retained pair, when it
is unique. Note that $\tau_0\in[\frac1{4R},U]$, since $b\le R$.

(a) The candidates are $\{i\}$, $\{1\}$ and, if $\{i,1\}$ is an edge, the pair.
By Weyl's inequality $\lambda_{\{i,1\}}(v)\ge\min(q_i-\delta,q_1)-\sqrt{w_{i1}}$,
so at $\tau_0$ the score of the pair exceeds the better singleton score by at
least $1-\tau_0R=\frac12$, and the pair never wins. The singleton $\{i\}$ has the
smaller score exactly when $q_i-\delta<q_1$, that is when $\delta>x_i$. Since
$|x_i|\le R$, both sides of $x_i$ lie in $[-3R,3R]$. Hence $\{1\}$ wins strictly
on $[-3R,x_i)$, $\{i\}$ wins strictly on $(x_i,3R]$, and the boundary is
$\delta=x_i$.

(b) Put $a=q_i-\bar\delta$, $c=q_j$ and $s=c-a=\bar\delta-\Delta_{ij}$, where
$\Delta_{ij}=q_i-q_j$ by (a). Then $R\le s\le3R$, and $\{i\}$ beats $\{j\}$ by
$\tau s>0$. By~\eqref{eq:lambda-sym} the pair has $\lambda_{\{i,j\}}(v)=a-g$, the
smaller eigenvalue of
$\bigl(\begin{smallmatrix}a&-\sqrt w\\-\sqrt w&c\end{smallmatrix}\bigr)$ with
$w=w_{ij}$, where
\[
 g=\frac{\sqrt{s^2+4w}-s}2=\frac{2w}{\sqrt{s^2+4w}+s}\in(0,\sqrt w\,] .
\]
So $\Phi(\{i,j\})-\Phi(\{i\})=1-\tau g$. Hence $\{i\}$ wins strictly below
$t_{ij}=1/g$, the pair wins strictly above it, and $1/g$ is the displayed
formula. Also $g\le\sqrt w\le R$ gives $t_{ij}\ge1/R>\frac1{4R}$. The formula
increases in $s$ and decreases in $w$, and $s\le3R$, $w\ge b^2$, $b\le R$, so
$t_{ij}\le(\sqrt{13}+3)R/(2b^2)<U/2$. At $\tau=t_{ij}$ the pair eigenvalue
$a-1/t_{ij}$ is a root of $(a-x)(c-x)=w$, which gives
$t_{ij}^{-1}(s+t_{ij}^{-1})=w$, the inversion formula.

(c) $B=S-q_1I$, since $S$ has diagonal $q_k$, entries $-\sqrt{w_{kl}}$ on edges
and zero off the graph. The positive unit vector $\sqrt\pi$ satisfies
$S\sqrt\pi=D_{\sqrt\pi}(-Q)\ones=0$. As $S\ge0$, the eigenvalue $0$ is
$\lambda_{\min}(S)$, and it is simple with eigenvector $\sqrt\pi$ by the theorem of
Perron and Frobenius for the irreducible nonnegative matrix $RI-S$. So
$\lambda=-q_1$, $S=B-\lambda I$, $\phi=\sqrt\pi$ and $\pi_i=\phi_i^2$, and
$Q=-D_\phi^{-1}SD_\phi$ is the definition of $S$. The coordinates used are $d-1$
values $x_i$ and $m$ values $t_{ij}$, and the fields and time scales of the design
depend only on $R$, $b$ and the labels.
\end{proof}

The reconstruction of a reversible generator from its exit rates and edge
products is classical. Kim and Landim~\cite[Lemma~5.1 and Corollary~5.3]{kl2024}
recover the same quantities from the values of $J_V$, and with the stationary law
from $J_V^{-1}(0)$ (their Lemma~5.2) deduce that $J_V$ determines a reversible
generator. Here the coordinates are read off the
boundaries of strict support regions of the weighted composition law, with a
design fixed in advance.

\begin{theorem}[minimality]\label{thm:minimality}
The reversible generators on the fixed connected graph $G$ form a real-analytic
manifold of dimension $m+d-1$, with global coordinates the edge products
$(w_{ij})$ and the exit differences $(q_i-q_1)$. Let $W$ be a nonempty open
subset of it, for example the set where $\|S\|_2<R$ and $\sqrt{w_{ij}}>b$ when
that set is nonempty. Let $k<m+d-1$, let $f\colon W\to\R^k$ be $C^1$ and let
$\ell\colon W\to\Lambda$ be locally constant with $\Lambda$ finite. Then the
observation $(f,\ell)$ is not injective on any nonempty open subset of $W$. So no
family of fewer than $m+d-1$ locally $C^1$ scalar observations with locally
constant finite labels is locally identifying.
\end{theorem}

\begin{proof}
Let $\mathcal Q_G$ be the set of reversible generators with $q_{ij}>0$
exactly on the edges of $G$, and $\kappa(Q)=((w_{ij})_E,(q_i-q_1)_{i\ne1})$.
Given $w\in(0,\infty)^m$ and $d\in\R^{d-1}$ (with $d_1=0$), form $B$ as in
Theorem~\ref{thm:reversible-reconstruction}(c), let $\lambda$ and $\phi$ be its least
eigenvalue and positive unit eigenvector, which is simple since $cI-B$ is
irreducible and nonnegative for large $c$, and put
$\Psi(w,d)=-D_\phi^{-1}(B-\lambda I)D_\phi$. This is a reversible generator on
$G$, with respect to $\phi^2$, with products $w_{ij}$ and exit rates
$d_i-\lambda$, so $\kappa\circ\Psi=\mathrm{id}$, and
Theorem~\ref{thm:reversible-reconstruction}(c), whose proof uses only
reversibility and irreducibility, gives $\Psi\circ\kappa=\mathrm{id}$. The map
$\Psi$ is real-analytic by the implicit function theorem at a simple
eigenvalue, and $\kappa$ is polynomial. So $\kappa$ is a global real-analytic
chart of $\mathcal Q_G$, of dimension $m+d-1$.

For the lower bound let $W'\subseteq W$ be nonempty and open. Shrinking it, we
may assume that $\ell$ is constant on $W'$ and that $W'$ is a ball in the chart
$\kappa$, so $f$ is a $C^1$ map on a ball of $\R^p$, $p=m+d-1$. Let $\rho$ be the
largest rank of $df$ on $W'$, attained at $x_0$. The rank is lower
semicontinuous, so it equals $\rho$ near $x_0$, and by the constant rank theorem
$f$ is constant on a $C^1$ submanifold through $x_0$ of dimension
$p-\rho\ge p-k\ge1$. So $(f,\ell)$ is not injective on $W'$.
\end{proof}

The chart is the continuous-time version of that of Wolfer and
Watanabe~\cite[Theorems~3, 5 and~6]{wolferwatanabe2021} for reversible kernels,
and the lower bound is the classical rank argument. Regularity is needed, since
interleaving the digits of the $m+d-1$ coordinates gives one injective scalar.

\begin{remark}[partial data]\label{rem:partial-data}
Suppose that the products are known on $J\subseteq E$ and the exit differences
along a comparison graph $K$ on $V$. In the chart $\kappa$ the generators with
these data are those with the given $w_e$ on $J$, any $w_e>0$ on $E\setminus J$,
and $d$ equal to the given one plus a vector that is constant on each component
of $K$ and zero on that of the state $1$. So they are parametrized by
$(0,\infty)^{m-|J|}\times\R^{c(K)-1}$, with $c(K)$ the number of components of
$K$, and the data identify $Q$ exactly when $J=E$ and $K$ is connected. For
instance, for $\ell>\sqrt2$ and $z=\ell-2/\ell$ the generators
\[
 Q_\ell=\begin{pmatrix}-\ell&2/\ell&z\\ \ell/2&-\ell&\ell/2\\ z&2/\ell&-\ell\end{pmatrix}
\]
satisfy detailed balance with $\nu=(1,4/\ell^2,1)$, have all exit rates equal
to $\ell$ and $w_{12}=w_{23}=1$, and have $w_{13}=z^2$. So the data of a spanning
tree miss a parameter.
\end{remark}

\begin{theorem}[stability]\label{thm:stability}
In the setting of Theorem~\ref{thm:reversible-reconstruction}, let the
observed dimensionless coordinates $d_i/R$ and $Rt_{ij}$ have sup-norm errors
at most $\varepsilon$, and let $\widehat Q$ be the reconstruction from the
observed values. There are $\varepsilon_0$ and $C$ depending only on $d$ and
$b/R$ such that for $\varepsilon\le\varepsilon_0$ the reconstruction
$\widehat Q$ is a reversible generator on $G$ and
\[
 \|\widehat Q-Q\|_2\le C\,R\,\varepsilon .
\]
Among the constants are the bound $\phi_{\min}\ge \phi_*=(b/R)^{d-1}/\sqrt d$
for the ground vector $\phi=\sqrt\pi$ and the bound
$\lambda_2(S)\ge g_*=2b\phi_*^2/(d-1)$ for the gap of $S$.
\end{theorem}

\begin{proof}
The reconstruction commutes with the scaling $Q\mapsto Q/R$, so we take $R=1$,
$\bar\delta=2$ and $\eta=b\le1$. Then $\sigma_{ij}=t_{ij}\ge1$ and
$a_{ij}=2-d_i+d_j\in[1,3]$, by Theorem~\ref{thm:reversible-reconstruction}. Hats
denote reconstructed values, and $\widehat d_1=d_1=0$.

\emph{Products.} Let $\varepsilon\le\eta^2/20$. Then $\widehat\sigma\ge\frac34$
and $|\widehat\sigma^{-1}-\sigma^{-1}|\le\frac43\varepsilon$, and
$w=a/\sigma+\sigma^{-2}$ gives
$|\widehat w-w|\le(\frac{20}3+\frac{28}9)\varepsilon<10\varepsilon$. So
$\widehat w\ge\eta^2/2$ and $|\sqrt{\widehat w}-\sqrt w|\le10\varepsilon/\eta$,
and the largest absolute row sum gives
$\|\widehat B-B\|_2\le\beta:=10d\varepsilon/\eta$. By Weyl's inequality
$|\widehat\lambda-\lambda|\le\beta$, so $\|\widehat S-S\|_2\le2\beta$. The matrix
$\widehat B$ has negative entries exactly on the edges of $G$, so the chart
$\Psi$ of the proof of Theorem~\ref{thm:minimality} makes $\widehat Q$ a
reversible generator with the graph $G$.

\emph{The constants.} Row $l$ of $S\phi=0$ reads
$q_l\phi_l=\sum_{k\sim l}\sqrt{w_{lk}}\phi_k$, so $\phi_l\ge\eta\phi_k$ on every
edge. Some $\phi_{k_0}$ is at least $1/\sqrt d$, and every state is joined to
$k_0$ by at most $d-1$ edges, so $\phi_{\min}\ge\phi_*$. For $x\perp\phi$ write
$x=\phi\circ f$. Then $\langle x,Sx\rangle=\sum_{\{k,l\}\in E}c_{kl}(f_k-f_l)^2$
with $c_{kl}=\sqrt{w_{kl}}\phi_k\phi_l\ge\eta\phi_*^2$, while
$\|x\|^2=\sum_k\pi_kf_k^2$ and $\sum_k\pi_kf_k=0$. In a spanning tree
$\mathcal T$ of $G$ the path from $k$ to $l$ has at most $d-1$ edges, so the
inequality of Cauchy and Schwarz gives
\[
 \|x\|^2=\tfrac12\sum_{k,l}\pi_k\pi_l(f_k-f_l)^2\le\tfrac{d-1}2\sum_{e\in\mathcal T}(\nabla_ef)^2
 \le\frac{d-1}{2\eta\phi_*^2}\langle x,Sx\rangle,
\]
and $\lambda_2(S)\ge g_*$.

\emph{The generator.} Write $\widehat\phi=c\phi+z$ with $z\perp\phi$ and $c>0$.
Since $Sz=(S-\widehat S)\widehat\phi$, we get $\|z\|\le2\beta/g_*$, and since
$1-c\le1-c^2=\|z\|^2$, also
$\gamma=\|\widehat\phi-\phi\|\le\sqrt2\|z\|\le2\sqrt2\beta/g_*$. This is the
bound of Davis and Kahan in its simplest form. Let $\varepsilon$ be so small
that $\gamma\le\phi_*/2$. Then $\widehat\phi_{\min}\ge\phi_*/2$,
$|\widehat\phi_l/\widehat\phi_k-\phi_l/\phi_k|\le4\gamma/\phi_*^2$ and
$\widehat\phi_l/\widehat\phi_k\le2/\phi_*$, so on an edge
\[
 |\widehat q_{kl}-q_{kl}|
 =\bigl|\sqrt{\widehat w_{kl}}\,\widehat\phi_l/\widehat\phi_k-\sqrt{w_{kl}}\,\phi_l/\phi_k\bigr|
 \le\frac{20\varepsilon}{\eta\phi_*}+\frac{4\gamma}{\phi_*^2}.
\]
The diagonal satisfies $|\widehat q_k-q_k|=|\widehat S_{kk}-S_{kk}|\le2\beta$, and
off the graph both matrices vanish. So every entry of $\widehat Q-Q$ is at most
$C_1(d,\eta)\varepsilon$, and $\|\widehat Q-Q\|_2\le dC_1\varepsilon$. The
conditions on $\varepsilon$ depend only on $d$ and $\eta$.
\end{proof}

The perturbation bounds are standard. The point is that the constants come from
$R$ and $b$ alone, with no separate assumption on the spectral gap.

\subsection{The non-reversible case}\label{sec:nonreversible-reconstruction}

Without reversibility the exit rates and the edge products no longer
determine $Q$, since the time reversal $Q^*=D_\pi^{-1}Q^\top D_\pi$ has the same
exit rates, the same products and the same stationary law. The right data are
the principal minors of $-Q$. As in Section~\ref{sec:spectral-geometry},
$\lambda_F(v)$ is the smallest real eigenvalue of $-Q_F-\operatorname{diag}v_F$.
For $A\subseteq F$ the principal submatrix of $-Q_F$ on $A$ is $-Q_A$, since a
killed generator keeps the full exit rates on its diagonal. So the principal
minors of $-Q_F$ are the numbers $\det(-Q_A)$, and they are the coefficients of
the multiaffine polynomial
\[
 R_F(w)=\det\bigl(-Q_F-\operatorname{diag}w\bigr)
 =\sum_{A\subseteq F}\det(-Q_A)\prod_{i\in F\setminus A}(-w_i),\qquad w\in\R^F,
\]
with $\det(-Q_\emptyset)=1$.

\begin{lemma}[irreducibility]\label{lem:RF-irreducible}
If $F$ is irreducible, then $R_F$ is irreducible over $\mathbb C$.
\end{lemma}

\begin{proof}
For $|F|=1$ the polynomial has degree $1$. Let $|F|\ge2$ and suppose that
$R_F=GH$ with $G$ and $H$ nonconstant. The polynomial $R_F$ has degree exactly
$1$ in each $w_i$, as its coefficient at $\prod_F(-w_i)$ is $1$, and degrees in $w_i$ add, so the sets $X$ and $Y$ of
variables of $G$ and $H$ partition $F$, both are nonempty, and $G$ and $H$ are
multiaffine. Scale them so that the coefficients of $\prod_X(-w_i)$ in $G$ and of
$\prod_Y(-w_i)$ in $H$ are $1$, and write
$G=\sum_{A\subseteq X}g_A\prod_{X\setminus A}(-w_i)$ and
$H=\sum_{B\subseteq Y}h_B\prod_{Y\setminus B}(-w_i)$. Comparing coefficients
gives $\det(-Q_{A\cup B})=g_Ah_B$, so $g_A=\det(-Q_A)$ and $h_B=\det(-Q_B)$ by
taking $B=\emptyset$ or $A=\emptyset$. Hence
\begin{equation}\label{eq:minor-split}
 \det(-Q_U)=\det(-Q_{U\cap X})\det(-Q_{U\cap Y})\qquad\text{for every }U\subseteq F.
\end{equation}

Since $G_F$ is strongly connected, it has an arc $x\to y$ with $x\in X$ and $y\in
Y$, and a shortest path from $y$ back to $x$ closes it into a simple directed
cycle that meets $X$ and $Y$. Let $U$ be a minimal set, for inclusion, among the
vertex sets of such cycles, and let $C$ be such a cycle on $U$. Expand $\det(-Q_U)$
over the permutations $\sigma$ of $U$. The permutations that preserve $U\cap X$
give $\det(-Q_{U\cap X})\det(-Q_{U\cap Y})$. Let $\sigma$ be another permutation with
a nonzero product. Every nontrivial cycle of $\sigma$ is a directed cycle of
$G_U$, and one of them meets $X$ and $Y$. By minimality its vertex set is $U$,
so $\sigma$ is a single cycle through all of $U$, and its term is
$\operatorname{sgn}\sigma\prod_{u\in U}(-q_{u\sigma(u)})=(-1)^{|U|-1}(-1)^{|U|}\prod_uq_{u\sigma(u)}
=-\prod_uq_{u\sigma(u)}$. Hence
\[
 \det(-Q_U)-\det(-Q_{U\cap X})\det(-Q_{U\cap Y})=-\sum_\sigma\prod_{u\in U}q_{u\sigma(u)}<0,
\]
the sum over the Hamiltonian cycles of $G_U$, which include $C$. This
contradicts~\eqref{eq:minor-split}.
\end{proof}

Al Ahmadieh~\cite[Lemma~5.1 and Corollary~5.6]{alahmadieh2023} proves over any
field that a matrix is irreducible exactly when $\det(\operatorname{diag}x+A)$
is an irreducible polynomial. The proof above is a short version for the
matrices $-Q_F$. The converse direction is immediate here, since a reducible
$-Q_F$ makes $R_F$ the product of the polynomials of its diagonal blocks.

\begin{lemma}[the Perron branch]\label{lem:perron-branch}
Let $F$ be irreducible, let $Q'$ be another generator whose face $F$ is
irreducible, and let $c\in\R$. Suppose that $\lambda'_F(v)=\lambda_F(v)+c$ for
all $v$ in a nonempty open set $\mathcal U\subseteq\R^F$. Then $-Q'_F$ and $-Q_F+cI$ have the
same principal minors, and $\lambda'_F=\lambda_F+c$ on all of $\R^F$. In
particular, for $c=0$, $-Q'_F$ and $-Q_F$ have the same principal minors.
\end{lemma}

\begin{proof}
Let $k=|F|$ and number the states of $F$ as $1,\dots,k$. Put
$\widetilde R(w)=\det(-Q_F+cI-\operatorname{diag}w)=R_F(w-c\ones)$, which is
irreducible over $\mathbb C$ by Lemma~\ref{lem:RF-irreducible}. For $k=1$ the
hypothesis says $q'_1=q_1+c$, which is the claim. Let $k\ge2$, and put
$\psi(v)=v+\lambda'_F(v)\ones$. For $v\in \mathcal U$ the number
$\lambda'_F(v)=\lambda_F(v)+c$ is an eigenvalue of $-Q'_F-\operatorname{diag}v$ and
of $-Q_F+cI-\operatorname{diag}v$, so $R'_F(\psi(v))=0=\widetilde R(\psi(v))$.

Both polynomials are affine in $w_1$. Write $\widetilde R=p_1w_1+p_0$ and
$R'_F=p'_1w_1+p'_0$ with $p_1,p_0,p'_1,p'_0$ polynomials in
$w'=(w_2,\dots,w_k)$. The coefficient of $\prod_{i\ge2}(-w_i)$ in $p_1$ is
$-1$, so $p_1\ne0$. The polynomial
$\mathcal D=p_1R'_F-p'_1\widetilde R=p_1p'_0-p'_1p_0$ does not
involve $w_1$ and vanishes on $\psi(\mathcal U)$. We claim that the projection of
$\psi(\mathcal U)$ to $w'$ contains a nonempty open set. Fix $v^0\in \mathcal U$ and consider
$y\mapsto(\psi(v^0+(0,y)))_{i\ge2}$ for small $y\in\R^{k-1}$. The Perron root of
$Q'_F+\operatorname{diag}v$ is simple, so $\lambda'_F$ is real-analytic with
$\partial\lambda'_F/\partial v_i=-\beta_i$, where $\beta=l\circ r>0$ for the
Perron vectors with $l\cdot r=1$, and $\sum_i\beta_i=1$. The derivative of the
map at $y=0$ is $I-\ones\beta'^\top$ with $\beta'=(\beta_i)_{i\ge2}$, of
determinant $\beta_1>0$, and the inverse function theorem proves the claim. So
$\mathcal D=0$ and $p_1R'_F=p'_1\widetilde R$. The irreducible $\widetilde R$
divides $p_1R'_F$, and it involves $w_1$, so it does not divide $p_1\ne0$, which
is free of $w_1$. Hence it divides $R'_F$. Both have total degree $k$ and the coefficient $1$ at
$\prod_i(-w_i)$, so $R'_F=\widetilde R$. Finally $\lambda'_F(v)$ is the smallest
real root of $t\mapsto R'_F(v+t\ones)=\widetilde R(v+t\ones)$, which is
$\lambda_F(v)+c$.
\end{proof}

Delmas, Dronnier and Zitt~\cite[Theorem~2.1]{ddz2024} prove a multiplicative
analogue, for the functions $\eta\mapsto\varrho(K\operatorname{diag}\eta)$ of
nonnegative matrices on $\{0,1\}^n$. The additive version above needs agreement
only on an open set, and uses the irreducibility of $R_F$ in place of their
induction over the cube. Cohen~\cite{cohen1981} proved that
$v\mapsto\varrho(A+\operatorname{diag}v)$ is convex for a Metzler matrix $A$,
which is the concavity of $\lambda_F$ in Proposition~\ref{prop:variational}, and
Friedland~\cite[Theorem~4.1]{friedland1981} derived it from a variational formula
of Donsker and Varadhan type and showed that for irreducible $A$ it is affine only
along the direction of $\ones$. These are statements about the map
$v\mapsto\varrho(A+\operatorname{diag}v)$. Lemma~\ref{lem:perron-branch} is the
inverse statement, that this map on an open set determines the principal minors
of $A$, and we did not find it there.

The next lemma replaces the Schur estimate of Lemma~\ref{lem:suppression},
which needs a symmetric matrix, by the bound of Collatz and Wielandt. For a
Metzler matrix $B$ and a positive vector $u$, the Perron root of $B$ is at most
$\max_i(Bu)_i/u_i$~\cite{bermanplemmons1994}.

\begin{lemma}[suppression without reversibility]\label{lem:suppression-nr}
Let $Q$ be a generator, $F$ a nonempty face, $\mathcal K\subset\R^F$ compact
and $\eta>0$. There is $\bar v_0$ such that for $\bar v\ge \bar v_0$, every field $v$ with
$v_F\in\mathcal K$ and $v_j=-\bar v$ off $F$, and every face $G\not\subseteq F$, we
have $\lambda_G(v)\ge\lambda_{G\cap F}(v)-\eta$ if $G\cap F\ne\emptyset$, and
$\lambda_G(v)\ge \bar v$ if $G\cap F=\emptyset$.
\end{lemma}

\begin{proof}
Put $P=G\cap F$, $S=G\setminus F$ and $B=Q_G+\operatorname{diag}v_G$, whose
Perron root is $-\lambda_G(v)$. If $P=\emptyset$, every row of $B$ has
off-diagonal sum at most $q_s$ and diagonal $-q_s-\bar v$, so its row sums are at most
$-\bar v$, and the bound of Collatz and Wielandt with $u=\ones$ gives $\lambda_G(v)\ge \bar v$.

Let $P\ne\emptyset$ and $\varrho_P(v)=-\lambda_P(v)$, and let $\mathbf J_P$ have
ones off the diagonal of $P$. For $\zeta>0$ the Perron root of the irreducible
Metzler matrix $Q_P+\operatorname{diag}v_P+\zeta\mathbf J_P$ is continuous in
$(\zeta,v_P)$ and decreases to $\varrho_P(v)$ as $\zeta\downarrow0$, so by Dini's
theorem some $\zeta$ makes it at most $\varrho_P(v)+\eta/2$ for all
$v_F\in\mathcal K$. Let $x$ be its Perron vector with $\max x=1$, and $x_*>0$ the
least entry of $x$ over $\mathcal K$. Test $B$ with $u=x$ on $P$ and $u=\xi$ on
$S$. For $p\in P$, since $\mathbf J_Px\ge0$,
\[
 (Bu)_p\le\bigl((Q_P+\operatorname{diag}v_P)x\bigr)_p+\xi q_{\max}
 \le(\varrho_P+\tfrac\eta2)x_p+\xi q_{\max}\le(\varrho_P+\eta)x_p
\]
once $\xi\le\eta x_*/(2q_{\max})$. For $s\in S$,
\[
 (Bu)_s\le\sum_{p\in P}q_{sp}x_p+\xi\Bigl(\sum_{s'\in S\setminus\{s\}}q_{ss'}-q_s-\bar v\Bigr)
 \le q_{\max}-\xi \bar v\le(\varrho_P+\eta)\xi
\]
once $\bar v\ge q_{\max}/\xi+|\varrho_P|+\eta$, which holds uniformly on
$\mathcal K$ for large $\bar v$. There are finitely many $P$ and $S$. So the Perron
root of $B$ is at most $\varrho_P(v)+\eta$, that is
$\lambda_G(v)\ge\lambda_P(v)-\eta$.
\end{proof}

\begin{lemma}[exposure by descent]\label{lem:exposure}
Let $Q$ be irreducible with a start of full support, and let $F$ be an
irreducible face with $|F|\ge2$. For $v_F\in\R^F$ and $\bar v>0$ write $(v_F,-\bar v)$
for the field equal to $v_F$ on $F$ and to $-\bar v$ off $F$. Outside a closed
nowhere dense set $\mathcal E\subset\R^F$, every $v^0$ has a unique partner
$G$, an irreducible face $G\subsetneq F$, with the following property. There are
an open neighbourhood $\mathcal U$ of $v^0$, a real-analytic function $\tau_1>0$ on $\mathcal U$,
and numbers $\delta>0$ and $\bar v_0$, such that for $\bar v\ge \bar v_0$ and the fields
$(v_F,-\bar v)$, on $O=\{(\tau,v_F):v_F\in \mathcal U,\ |\tau-\tau_1(v_F)|<\delta\}$,
\begin{enumerate}[(i)]
\item the points with $\tau>\tau_1(v_F)$ lie in the strict winning region of $F$,
and those with $\tau<\tau_1(v_F)$ in that of $G$,
\item on the hypersurface $\tau=\tau_1(v_F)$ we have
$\lambda_F(v)=\lambda_G(v)-(|F|-|G|)/\tau_1(v)$.
\end{enumerate}
The neighbourhood $\mathcal U$, the function $\tau_1$ and the numbers $\delta$
and $\bar v_0$ depend on $Q$.
For $F=V$ there is nothing to suppress. Hence the controlled boundaries
determine, for every irreducible face $F$, the function $\lambda_F$ on all of
$\R^F$, up to one additive constant common to all faces, and
$\lambda_V(0)=0$ fixes the constant. In this comparison of $Q$ with a second
generator $Q'$ the ranges of $\tau$, $\delta$ and $\bar v$ are chosen for both,
so they depend on $Q'$, and unlike in the reversible case no design is fixed in
advance. The statement is one of identifiability.
\end{lemma}

\begin{proof}
Let $k=|F|$. For an irreducible face $H\subsetneq F$, singletons included, put
\[
 \rho_H(v)=\frac{k-|H|}{\lambda_H(v)-\lambda_F(v)},\qquad v\in\R^F .
\]
The denominator is positive by Lemma~\ref{lem:face-monotone}, and $\rho_H$ is
real-analytic, since the Perron roots of irreducible Metzler matrices are
simple. Note that $\Phi_{\tau,v}(F)<\Phi_{\tau,v}(H)$ exactly when $\tau>\rho_H(v)$.

\emph{Generic partners.} Let $H_1\ne H_2$ be two such faces with some $j\in
H_1\setminus H_2$, and take $v=te_j$ with $t\to\infty$. Then $\lambda_{H_2}$
stays fixed, $\lambda_F\le\lambda_{H_1}\le q_j-t$ since a Perron root is at
least every diagonal entry, and $\lambda_F\ge-t$ by the bound of Collatz and
Wielandt with $u=\ones$. Hence $\rho_{H_2}\to0$, while
$\rho_{H_1}\ge(k-|H_1|)/q_j>0$, as $q_j>0$ on the irreducible face $F$. So $\rho_{H_1}-\rho_{H_2}$ is a nonzero
real-analytic function, with a closed nowhere dense zero set. The set
$\mathcal E$ where the maximum of $\rho_H$ over $H$ is attained twice is closed
and lies in the finite union of these zero sets.

\emph{The boundary.} Let $v^0\notin\mathcal E$, let $G$ be the unique maximizer of
$\rho_H(v^0)$ and put $\tau_1=\rho_G$. Choose $\delta>0$ with
$3\delta<\tau_1(v^0)$ and $\rho_H(v^0)<\tau_1(v^0)-3\delta$ for every irreducible
$H\subsetneq F$ other than $G$. By continuity there is a neighbourhood $\mathcal U$ of
$v^0$ with compact closure on which $\tau_1>2\delta$ and $\rho_H<\tau_1-2\delta$
for all these $H$. Then $\tau-\rho_H\ge\delta$ on the closure of $O$, and
$\Phi(H)-\Phi(F)=(\lambda_H-\lambda_F)(\tau-\rho_H)\ge\kappa$ for some $\kappa>0$.
Also $\Phi(F)-\Phi(G)=(\lambda_G-\lambda_F)(\tau_1(v)-\tau)$ is negative above the
hypersurface and positive below it, and a reducible face loses to one of its
components (Proposition~\ref{prop:controlled-score}). So among the faces inside
$F$ the least score is attained only by $F$ above and only by $G$ below. Apply
Lemma~\ref{lem:suppression-nr} to the closure of $\mathcal U$, with
$\eta=1/(2\tau_+)$ for an upper bound $\tau_+$ of $\tau$ on $O$. Let
$G'\not\subseteq F$ be admissible. If $P=G'\cap F$ is nonempty and $C$ is a
component of $P$ with $\lambda_C(v)=\lambda_P(v)$, then
$\Phi(G')\ge|P|+\tau\lambda_P(v)-\frac12\ge\Phi(C)+\frac12$. If $P$ is empty,
$\Phi(G')\ge\tau_-\bar v$ for a lower bound $\tau_->0$ of $\tau$ on $O$, which
exceeds the scores inside $F$ by $\frac12$ for large $\bar v$. This proves (i),
and the tie $\Phi(F)=\Phi(G)$ on the hypersurface is (ii).

\emph{The induction.} Let $Q'$ be a generator with the same controlled
boundaries. By Lemma~\ref{lem:face-monotone} the face $V$ is the strict winner
for $Q$ on an open set of points $(\tau,v)$, so it wins strictly for $Q'$ at some
points, and Proposition~\ref{prop:controlled-score} makes $Q'$ irreducible. For
$i\ne1$ take $F=\{1,i\}$ for both generators at once, with the field $\delta$ on
$i$ and $0$ on $1$, $\tau$ in a compact subinterval of
$(0,1/(2\max(q_{\max},q'_{\max}))]$, $\delta$ in a compact interval containing
$q_i-q_1$ and $q'_i-q'_1$ in its interior, and $\bar v$ above the thresholds of
Lemma~\ref{lem:suppression-nr} for $Q$ and $Q'$. An irreducible pair has
eigenvalue at least $\min(q_1,q_i-\delta)-\sqrt{w_{1i}}$ and loses to the better
singleton by at least $\frac12$, for both generators, and
Lemma~\ref{lem:suppression-nr} removes the other faces. So the boundaries
$\delta=q_i-q_1$ and $\delta=q'_i-q'_1$ agree. Put $c=q'_1-q_1$, so that
$\lambda'_{\{i\}}=\lambda_{\{i\}}+c$ for every $i$. We show by induction on $|F|$ that every irreducible face $F$ of
$Q$ is irreducible for $Q'$ and that $\lambda'_F=\lambda_F+c$ on $\R^F$. Let
$|F|\ge2$, choose $v^0\notin\mathcal E$ with partner $G$, and take $O$ as above.
Outside a closed nowhere dense subset of $O$, the faces $F$ above the
hypersurface and $G$ below it are strict winners for $Q'$ too, hence irreducible
for $Q'$. By continuity of the scores of $Q'$, $\Phi'(F)=\Phi'(G)$ on the
hypersurface. With (ii) and the induction hypothesis for $G$,
\[
 \lambda'_F(v)=\lambda'_G(v)-\frac{|F|-|G|}{\tau_1(v)}=\lambda_G(v)+c-\frac{|F|-|G|}{\tau_1(v)}=\lambda_F(v)+c
 \qquad(v\in \mathcal U).
\]
Lemma~\ref{lem:perron-branch} extends this to all of $\R^F$. For $F=V$ both
generators have $\lambda_V(0)=\lambda'_V(0)=0$, since $\ones$ is a positive null
vector of $Q$ and of $Q'$. So $c=0$.
\end{proof}

The induction is an algorithm. For each face choose a field at random,
decrease $\tau$ from a large value, and read the first change of label and its
position. The lemma is the place where the boundaries of all faces, and not only
those of pairs, are used.

For a face $F$ that is irreducible for $Q$ and a field $v$, let $\alpha^*_F(v)$
be the minimizer of $J_F-v\cdot\alpha$ over $\Delta_F$
(Proposition~\ref{prop:variational}) and $K_F(v)=C_F(\alpha^*_F(v))$, the
constant of Definition~\ref{def:tilt} at that point. So $K_F(0)=K_F$, and
$K_F(v)$ is the constant of the face maxima with a field
(Theorem~\ref{thm:facemax}).

\begin{theorem}[non-reversible reconstruction]\label{thm:nonreversible}
Let $Q$ be irreducible with a start of full support, and suppose that the
controlled boundaries of all irreducible faces are observed.
\begin{enumerate}[(a)]
\item The boundaries determine every principal minor of $-Q$.
\item Every generator $Q'$ with the same principal minors is irreducible, has
the same spectrum as $Q$, and the set of such generators is finite. It
contains $Q$ and its time reversal $Q^*$.
\item If $Q$ has no cut, that is, for every partition $V=X\cup Y$ with
$|X|,|Y|\ge2$ one of the blocks $Q[X,Y]$, $Q[Y,X]$ has rank at least $2$, then
the set in (b) is $\{Q,Q^*\}$, which is a single point exactly when $Q$ is
reversible.
\item In general the set in (b) is the set of generator points of the
cut-transpose orbit of $-Q$. Each cut-transpose of an irreducible Z-matrix $C$
across a rank-one cut $X$, with $Y=V\setminus X$,
\[
 C=\begin{pmatrix}C_{XX}&-ab^\top\\-cu^\top&C_{YY}\end{pmatrix}\mapsto
 \begin{pmatrix}C_{XX}&-ac^\top\\-bu^\top&C_{YY}^\top\end{pmatrix},\qquad a,b,c,u\ge0,
\]
is again an irreducible Z-matrix with the same principal minors, and exactly
one positive diagonal similarity turns it into a generator. Fewer than $2d$
cut-transposes and one optional transposition reach every point of the fiber.
If $Q$ is reversible, the fiber is $\{Q\}$.
\item If $d=3$, the fiber is $\{Q,Q^*\}$ with no condition.
\item There is a finite union $\mathcal Z$ of proper linear subspaces of $\R^V$
such that, for every start $\mu$ of full support outside $\mathcal Z$, every
member $Q'\ne Q$ of the fiber has $K'_V(v)\ne K_V(v)$ for all $v$ in a dense open
set of fields. If $Q$ is not reversible, then $\pi\in\mathcal Z$, and with the
start $\pi$ the law of the composition is the same under $Q$ and $Q^*$ for every
$N$ and $T$ with $Tq_{\max}\le N$. If $Q'$ is the generator point of one
cut-transpose across $X$, then
for every start $K'_F=K_F$ on every face $F\subseteq X$ and $K'_F=K^*_F$ on every
face $F\subseteq Y$.
\end{enumerate}
\end{theorem}

Part (f) is stated for the amplitudes. To read $K_F(v)$ off the boundary
positions one needs the crossing equation of Theorem~\ref{thm:crossing} for the
weighted law, which we do not prove here.

\emph{What is borrowed.} The principal-minor fiber of an
irreducible matrix is known, up to diagonal equivalence and transposition.
Loewy~\cite[Theorem~1]{loewy1986} identifies a matrix of order $n\ge4$ without a
cut from its principal minors, and Hartfiel and
Loewy~\cite[Theorem~3]{hartfielloewy1984} do the same for $n\le3$ with no
condition. Chatterjee, Ghosh, Gurjar and Raj combine the two in their
Lemma~2.8~\cite{cggr2024}. For $n\ge4$, Al
Ahmadieh~\cite[Theorem~4.1 and Lemma~4.5]{alahmadieh2023} shows that the fiber is
a single class exactly when there is no cut or the matrix is diagonally
equivalent to a symmetric one, and Chatterjee, Ghosh, Gurjar and
Raj~\cite[Theorem~1.1]{cggr2024} describe the whole fiber by cut-transposes.
Kim and Landim~\cite[Section~5]{kl2024} extract the minors of orders one and two
from $J_V$ (their Lemma~5.1), deduce with the stationary law that $J_V$
determines a reversible generator (their Corollary~5.3), and note that the
uniform $3$-cycle and its reversal share $J_V$ (their Example~5.4).
Corollary~\ref{cor:dv-fiber} answers the question left open by that example. In
contrast, their Corollary~5.8 shows that the joint functional of the empirical
measure and the empirical current determines the generator when every state is
recurrent. So parts (c) and (e) and the reversible case of (d) are corollaries of
the matrix theorems together with the normalization (N) in the proof below. What
the paper adds is listed in Section~\ref{sec:intro-reconstruction}.

\begin{proof}
We first record a normalization, which (b) to (e) use.

\emph{(N) The generator point.} Let $C$ be an irreducible Z-matrix with the
same principal minors as $-Q$. The characteristic polynomial
$\det(tI-C)=\sum_{s=0}^d(-1)^st^{d-s}E_s(C)$, with $E_s$ the sum of the principal
minors of order $s$, is that of $-Q$. So $-C$ has the spectrum of $Q$, whose
eigenvalue of largest real part is the simple eigenvalue $0$. The theorem of
Perron and Frobenius for the irreducible Metzler matrix $-C$ gives a vector
$x>0$, unique up to a factor, with $Cx=0$. With $D=\operatorname{diag}x$ the
matrix $-D^{-1}CD$ has off-diagonal entries $-c_{ij}x_j/x_i\ge0$ and zero row
sums, so it is an irreducible generator. If $D'$ is real, diagonal and
invertible and $-D'^{-1}CD'$ has zero row sums, then $C(D'\ones)=0$, so $D'$ is a
multiple of $D$ and gives the same matrix. We call $-D^{-1}CD$ the generator
point of $C$. For $C=-Q$ it is $Q$, and for $C=-Q^\top$, whose null vector is
$\pi$, it is $Q^*$.

\emph{(R) Real similarities.} If $-Q'=D(-Q)D^{-1}$ or $-Q'=D(-Q)^\top D^{-1}$ for
a complex invertible diagonal $D$, then $D$ may be taken real. In the first case
$q'_{ij}=d_iq_{ij}/d_j$, so $d_i/d_j=q'_{ij}/q_{ij}$ is real on every edge of the
graph of $Q$. That graph is strongly connected, so every ratio $d_i/d_1$ is a
product of such ratios and is real, and $D/d_1$ is a real diagonal similarity.
The second case is the same with $q_{ji}$ in place of $q_{ij}$.

(a) Let $Q'$ be a generator with the same controlled boundaries. The induction
in the proof of Lemma~\ref{lem:exposure} shows that $Q'$ is irreducible and gives
$\lambda'_V=\lambda_V$ on $\R^V$, and Lemma~\ref{lem:perron-branch} with $F=V$ and
$c=0$ gives equal principal minors.

(b) Let $Q'$ be a generator with the same principal minors, and suppose that
$-Q'$ is reducible. Then there is a proper nonempty $C\subset V$ with
$Q'[V\setminus C,C]=0$, and $\det(-Q')=\det(-Q'_C)\det(-Q'_{V\setminus C})$. Both
factors are proper principal minors, so they equal $\det(-Q_C)$ and
$\det(-Q_{V\setminus C})$. By Lemma~\ref{lem:face-monotone} every irreducible
component of a proper face has a positive exit rate, so every eigenvalue of
$-Q_C$ has positive real part and $\det(-Q_C)>0$, and likewise
$\det(-Q_{V\setminus C})>0$. This contradicts $\det(-Q')=\det(-Q)=0$. So $Q'$ is
irreducible, and the same argument applies to every real matrix with the
principal minors of $-Q$. This is the case $A=-Q$ of Al
Ahmadieh~\cite[Corollary~5.6]{alahmadieh2023}, which states it for $n\ge4$ over
any field. The spectra agree by the expansion in (N). The time
reversal $Q^*$ is a generator, since $Q^*\ones=D_\pi^{-1}Q^\top\pi=0$ and its
off-diagonal entries are $\pi_jq_{ji}/\pi_i\ge0$, and
$-Q^*=D_\pi^{-1}(-Q)^\top D_\pi$ has the principal minors of $-Q$, since
transposition and diagonal similarity preserve them. Finiteness follows from (d).

(c) For $d=2$ a generator is fixed by $q_1$ and $q_2$, which are principal
minors, and for $d=3$ the claim is (e). Let $d\ge4$. Loewy's
theorem~\cite[Theorem~1]{loewy1986} applies to the irreducible matrix $-Q$
without a cut, with no condition on $-Q'$. It gives an invertible diagonal
$D$ with $-Q'=D(-Q)D^{-1}$ or $-Q'^\top=D(-Q)D^{-1}$, and by (R) we may take $D$
real. In the
first case (N) with $C=-Q$ gives $Q'=Q$, and in the second case
$Q'=D^{-1}Q^\top D$ and (N) with $C=-Q^\top$ gives $Q'=Q^*$. Finally $Q^*=Q$
exactly when $\pi_iq_{ij}=\pi_jq_{ji}$ for all $i,j$, that is when $Q$ is
reversible.

(d) Let $C$ be an irreducible Z-matrix with the principal minors of $-Q$, and
let $X$ be a cut. The block $C[X,Y]$ is nonzero, by irreducibility, and has
nonpositive entries and rank $1$. Every row is a nonnegative multiple of a fixed
nonzero row, so $C[X,Y]=-ab^\top$ with $a,b\ge0$, and likewise
$C[Y,X]=-cu^\top$. The cut-transpose $C'$ is then a Z-matrix. It has the same
principal minors. Indeed, for $U\subseteq V$ write $U_X=U\cap X$ and $U_Y=U\cap
Y$. If $U_X$ or $U_Y$ is empty, both minors agree since a determinant is
invariant under transposition. Otherwise replace $C_{XX}$ by $C_{XX}+tI$ with
$t$ large, so that the block of $U_X$ is invertible. The Schur complement gives
\[
 \det C_U=\det C_{U_X}\det\bigl(C_{U_Y}-\kappa\,c_{U_Y}b_{U_Y}^\top\bigr),\qquad
 \det C'_U=\det C_{U_X}\det\bigl(C_{U_Y}^\top-\kappa\,b_{U_Y}c_{U_Y}^\top\bigr),
\]
with the same scalar $\kappa=u_{U_X}^\top C_{U_X}^{-1}a_{U_X}$. The two right sides
are equal by transposition. Both sides are polynomials in $t$ that agree for
large $t$, so they agree at $t=0$. By (b), $C'$ is irreducible, and by (N) it
has exactly one generator point.

Now let $Q'$ be a generator in the fiber, irreducible by (b). The theorem of
Chatterjee, Ghosh, Gurjar and Raj~\cite[Theorem~1.1]{cggr2024} gives a sequence
$-Q=C_0,C_1,\dots,C_k$ with $k<2d$, each $C_{s+1}$ a cut-transpose of $C_s$
across a cut of $C_s$~\cite[Definition~2.10]{cggr2024}, such that $-Q'$ is
diagonally similar to $C_k$ or to $C_k^\top$. Their cut-transpose uses a
particular rank-one factorization, and another choice changes it by a real
diagonal similarity $\operatorname{diag}(\alpha I_X,I_Y)$~\cite[Remark~2.11]{cggr2024}.
A real invertible diagonal similarity of $C$ by $\operatorname{diag}(D_X,D_Y)$
changes the cut-transpose by $\operatorname{diag}(D_X,D_Y^{-1})$, as a short
computation with the blocks shows, and does not change the cuts. So, up to real
invertible diagonal similarity, which (N) accepts, every $C_s$ is obtained by our
cut-transposes, and $Q'$ is the generator point of $C_k$ or of $C_k^\top$.
Conversely every such generator point lies in the fiber. There are fewer than
$2^d$ cuts at each step and fewer than $2d$ steps, so the fiber is finite, which
completes (b).

If $Q$ is reversible, $S=D_{\sqrt\pi}(-Q)D_{\sqrt\pi}^{-1}$ is an irreducible
symmetric Z-matrix, and a symmetric matrix is its own cut-transpose. Indeed, if
$S[X,Y]=-ab^\top$, then $S[Y,X]=-ba^\top$, so we may take $c=b$ and $u=a$. By the
rule for diagonal similarities above, the cut-transpose of a matrix diagonally
similar to $S$ is again diagonally similar to $S$, and so is the transpose
$D^{-1}SD$ of $DSD^{-1}$. Since $C_0=-Q$ is diagonally similar to $S$, so is every
$C_s$. Hence $-Q'$ is diagonally similar to $-Q$, and (N) with $C=-Q$ gives
$Q'=Q$. See also Engel and Schneider~\cite[Theorem~3.5]{engelschneider1980} and Al
Ahmadieh~\cite[Lemma~6.1 and Theorem~6.2]{alahmadieh2023}.

(e) Let $d=3$ and let $Q'$ be in the fiber, so $Q'$ is irreducible by (b). The
theorem of Hartfiel and Loewy~\cite[Theorem~3]{hartfielloewy1984}, as restated
in~\cite[Lemma~2.8]{cggr2024}, says that for an irreducible matrix $A$ of order $2$
or $3$ over a field, every matrix with the principal minors of $A$ is diagonally
similar to $A$ or to $A^\top$ over that field. So $-Q'=D(-Q)D^{-1}$ or
$-Q'=D(-Q)^\top D^{-1}$ with $D$ diagonal and invertible, real by (R), and (N) gives
$Q'=Q$ or $Q'=Q^*$, as in (c).

(f) Let $Q'$ be in the fiber and let $F$ be irreducible for $Q$. The matrices $-Q_F$
and $-Q'_F$ have the same principal minors, so $R'_F=R_F$, which is irreducible by
Lemma~\ref{lem:RF-irreducible}. A reducible $-Q'_F$ would make $R'_F$ the product of
the polynomials of its diagonal blocks, so $F$ is irreducible for $Q'$. Then
$\lambda'_F(v)$ and $\lambda_F(v)$ are the smallest real roots of the same
polynomial $t\mapsto R_F(v+t\ones)$, so $\lambda'_F=\lambda_F$, and $J'_F=J_F$ on
$\Delta_F^\circ$ by the Legendre duality of Proposition~\ref{prop:variational}.
So $\alpha^*_F(v)$ and, by Proposition~\ref{prop:constant}(a), $\det\Sigma$ at that
point agree for $Q$ and $Q'$. Let $r$ and $l$ be the right and left Perron vectors
of $Q_F+\operatorname{diag}v_F$ with $l\cdot r=1$, so that $\alpha^*_F(v)=l\circ r$.
The computation after Definition~\ref{def:tilt}, with $q-v$ in place of $q$, shows
that $r$ attains the supremum defining $J_F(\alpha^*_F(v))$. So
$r_{\alpha^*_F(v)}$ is a multiple of $r$, $l_{\alpha^*_F(v)}$ the matching multiple
of $l$, and
\[
 K_F(v)=\frac{(\mu_F\cdot r)(l\cdot\ones)}{(2\pi)^{(k-1)/2}\sqrt{\det\Sigma_{\alpha^*_F(v)}}}.
\]
Hence $K'_F(v)/K_F(v)=a'_F(v)/a_F(v)$ with $a_F(v)=(\mu_F\cdot r)(l\cdot\ones)$.

Take $F=V$ and put $g(v)=r(v)\,(l(v)\cdot\ones)$, which does not depend on the
scale of $r$, so that $a_V(v)=\mu\cdot g(v)$. The map $g$ is real-analytic on
$\R^V$. Let $W$ be the span of the vectors $g'(v)-g(v)$ over $v\in\R^V$. If
$\mu\notin W^\perp$, the analytic function $a'_V-a_V=\mu\cdot(g'-g)$ is not
identically zero, so it is nonzero on a dense open set. We show that $W\ne0$.
Suppose $g'=g$. Then $r'(v)$ is a multiple of $r(v)$, and both matrices have the
eigenvalue $\lambda_V(v)$, so $(Q'-Q)r(v)=0$ for every $v$. Take $v=te_j$ with
$t\to\infty$. Then $\lambda_V\le q_j-t\to-\infty$, and row $i\ne j$ of the
eigenvalue equation gives $(q_i-\lambda_V)r_i\le q_{\max}\max r$, so
$r(te_j)/\max r\to e_j$. Hence $(Q'-Q)e_j=0$ for every $j$ and $Q'=Q$. So for
$Q'\ne Q$ the space $W^\perp$ is proper, and $\mathcal Z$ is the union of these
spaces over the finitely many $Q'\ne Q$ in the fiber.

For the time reversal, $-Q^*-\operatorname{diag}v=D_\pi^{-1}(-Q-\operatorname{diag}v)^\top
D_\pi$, so $r^*=D_\pi^{-1}l$ and $l^*=D_\pi r$, with $l^*\cdot r^*=1$. At
$\mu=\pi$ this gives $a^*_V(v)=(\ones\cdot l)(\pi\cdot r)=a_V(v)$ for all $v$, so
$\pi\in\mathcal Z$ when $Q^*\ne Q$. More is true. The matrix $P=I+hQ$ has
stationary law $\pi$, and $P^*=D_\pi^{-1}P^\top D_\pi=I+hQ^*$ is stochastic for
$hq_{\max}\le1$. For a path $x_1\cdots x_N$ we have
$\pi_{x_1}\prod_tP(x_t,x_{t+1})=\pi_{x_N}\prod_tP^*(x_{t+1},x_t)$. Reversal is a
bijection of the paths that keeps the composition, so the law of the composition
under $Q$ from $\pi$ equals that under $Q^*$ from $\pi$, for every $N$ and $T$
with $hq_{\max}\le1$.

Finally let $Q'$ be the generator point of the cut-transpose $C'$ of $C=-Q$ across
$X$, and put $s_0=(b\cdot\ones)/(c\cdot\pi_Y)$. The rows $X$ of $Q\ones=0$ give
$-Q_X\ones=(b\cdot\ones)a$, the columns $Y$ of $\pi^\top Q=0$ give
$-Q_Y^\top\pi_Y=(\pi_X\cdot a)b$, and summing the columns $X$ of $\pi^\top Q=0$
gives $(\pi_X\cdot a)(b\cdot\ones)=(\pi_Y\cdot c)(u\cdot\ones)$. With these
identities, $x=(\ones_X,s_0\pi_Y)$ satisfies $C'x=0$. So
$Q'=-D_x^{-1}C'D_x$, and its killed blocks are $Q'_X=Q_X$ and
$Q'_Y=D_{\pi_Y}^{-1}Q_Y^\top D_{\pi_Y}=Q^*_Y$. The constant $K_F(v)$ of a face
depends only on $Q_F$ and $\mu_F$, which gives the last claim.
\end{proof}

A cut-transpose need not be $Q$ or $Q^*$. Take $d=4$, $X=\{1,2\}$ and
\[
 Q=\begin{pmatrix}-4.7&0.7&1&3\\1.3&-9.3&2&6\\2&2&-4.4&0.4\\1&1&2&-4\end{pmatrix},
\]
whose blocks $Q[X,Y]=(1,2)^\top(1,3)$ and $Q[Y,X]=(2,1)^\top(1,1)$ have rank one.
To four decimals, the generator point of the cut-transpose of $-Q$ across $X$ is
\[
 Q'=\begin{pmatrix}-4.7&0.7&2.3810&1.6190\\1.3&-9.3&4.7619&3.2381\\
 0.84&0.84&-4.4&2.72\\1.8529&1.8529&0.2941&-4\end{pmatrix}.
\]
Its entries are rational, for instance $q'_{13}=50/21$, $q'_{41}=63/34$ and
$q'_{43}=5/17$. It has the fifteen principal minors of $-Q$, with $Q'_X=Q_X$ and
$Q'_Y=Q^*_Y$, but $Q'\ne Q^*$, since $q^*_{12}=351/530\approx0.6623$.

\begin{remark}[three states by hand]\label{rem:three-states}
For $d=3$ the minors of orders $1$ and $2$ give the exit rates and the products
$w_{ij}$, and $\det(-Q)=0$ gives the sum and the product $w_{12}w_{23}w_{13}$ of
the two oriented $3$-cycle products $c=q_{12}q_{23}q_{31}$ and
$c'=q_{13}q_{32}q_{21}$. So $Q'$ has the cycle products of $Q$ or of $Q^*$
along every simple cycle, hence differs from it by a potential, and (N) gives
(e) directly. The uniform $3$-cycle and its reversal, Example~5.4 of Kim and
Landim~\cite{kl2024}, are the case $c=1$, $c'=0$.
\end{remark}

\begin{proof}[Proof of Corollary~\ref{cor:dv-fiber}]
(i) implies (iii). If $J_V=J'_V$ on $\Delta_V$, the variational formula
(Proposition~\ref{prop:variational}) gives
$\lambda_V(v)=\min_{\Delta_V}(J_V-v\cdot\alpha)=\lambda'_V(v)$ for every $v$, and
Lemma~\ref{lem:perron-branch} with $\mathcal U=\R^V$ and $c=0$ gives equal principal
minors.

(iii) implies (i). Equal principal minors give $R'_V=R_V$. For every $v$,
$\lambda_V(v)$ is the smallest real root of $t\mapsto R_V(v+t\ones)$, so
$\lambda'_V=\lambda_V$ on $\R^V$. By Proposition~\ref{prop:variational}, $J_V$ is
the Legendre transform of $v\mapsto-\lambda_V(v)$ on $\Delta_V^\circ$, so
$J'_V=J_V$ there. On the boundary we use that $J_V$ is convex and lower
semicontinuous, as a supremum of affine functions of $\alpha$. For
$\alpha\in\Delta_V$ and $\beta\in\Delta_V^\circ$ the points
$\alpha_t=(1-t)\alpha+t\beta$ with $t\in(0,1]$ are interior. Convexity gives
$\limsup_{t\downarrow0}J_V(\alpha_t)\le J_V(\alpha)$, and lower semicontinuity the
reverse inequality. So $J_V(\alpha)=\lim_{t\downarrow0}J_V(\alpha_t)$, likewise
for $J'_V$, and $J'_V(\alpha)=J_V(\alpha)$.

(ii) implies (iii). This is Theorem~\ref{thm:nonreversible}(a).

(iii) implies (ii). Let $F$ be a face. Equal principal minors give $R'_F=R_F$.
By Lemma~\ref{lem:RF-irreducible} and the remark after it, $F$ is irreducible for
$Q$ exactly when $R_F$ is irreducible, so $Q$ and $Q'$ have the same irreducible
faces. On such a face $\lambda'_F(v)$ and $\lambda_F(v)$ are the smallest real
roots of $t\mapsto R_F(v+t\ones)$, so they agree, and the scores
$\Phi_{\tau,v}$ of the irreducible faces agree. With a start of full support
the irreducible faces are admissible, and every reducible admissible face has a
larger score than one of its components (Proposition~\ref{prop:controlled-score}).
So for both generators the strict winner at $(\tau,v)$ is the unique minimizer
of $\Phi_{\tau,v}$ over the irreducible faces, when it exists, and the strict
winners agree at every point.
\end{proof}

\subsection{Finite horizons}\label{sec:finite-horizon}

Theorem~\ref{thm:reversible-reconstruction} reads the coordinates off the
limiting boundaries. At a finite horizon the observed boundaries are shifted,
and the shift has an explicit form, uniformly over a class of chains and
starts. Throughout, the connected graph $G$ on $[d]$ and the numbers $0<b\le R$
and $0<\mu_*\le1/d$ are fixed. Let $\mathcal C$ be the set of reversible
irreducible generators on $G$ with $\|S\|_2\le R$ and $\sqrt{w_{ij}}\ge b$ on
every edge, and let $\mathcal M$ be the set of starts with
$\min_i\mu_i\ge\mu_*$. For a pair $(i,j)$ and $\delta\in\R$ let $v^{ij}(\delta)$
be the design field of Section~\ref{sec:reversible-reconstruction}, and let
$\mathcal V$ be the compact set of these fields with $|\delta|\le3R$. For a
face $F$ and a field $v$ put
\[
 M_F(T,v)=\max\bigl\{P_N(n)\,e^{Tv\cdot n/N}:\supp n=F\bigr\},
\]
with $\max\emptyset=0$, and recall $\alpha^*_F(v)$ and $K_F(v)$ from before
Theorem~\ref{thm:nonreversible}.

\begin{proposition}[uniform face maxima on the class]\label{prop:uniform-facemax}
Let $G$ induce a connected graph on the face $F$, with $k=|F|\ge2$. There are
$N_0$ and $C$, depending only on $G$, $R$, $b$ and $\mu_*$, such that for all
$Q\in\mathcal C$, $\mu\in\mathcal M$ and $v\in\mathcal V$ the following hold.
\begin{enumerate}[(a)]
\item For all $N\ge2$ and $T>0$ with $TR\le N/2$,
$M_F(T,v)\le C(T/N)^{k-1}e^{-T\lambda_F(v)+CT^2/N}$.
\item For $N\ge N_0$ and $1\le T\le N^{1/3}$, we have $C^{-1}\le K_F(v)\le C$ and
\[
 M_F(T,v)=K_F(v)\,N^{-(k-1)}T^{(k-1)/2}e^{-T\lambda_F(v)}(1+E),\qquad
 |E|\le C\Bigl(\frac1T+\frac{T^2}N\Bigr).
\]
\end{enumerate}
For a singleton, $M_{\{i\}}(T,v)=\mu_i(1-Tq_i/N)^{N-1}e^{Tv_i}$ exactly.
\end{proposition}

\begin{proof}
For fixed $(Q,\mu)$, part (b) is Theorem~\ref{thm:facemax} with the compact
set $\mathcal V$ of fields, whose constants already do not depend on
$v\in\mathcal V$. Part (a) is~\eqref{eq:upper-field}, that is, Lemma~\ref{lem:upper}
with the Perron tilt $\theta=q-v-\lambda_F(v)\ones\in V_F$, and the singleton is
the path $i^N$. It remains to show that no constant depends on $(Q,\mu)$ either.

\emph{The parameter set is compact.} Let $\mathcal P=\mathcal C\times\mathcal M\times\mathcal V$.
On an edge $q_{ij}\ge b^2/R$, since $q_{ij}q_{ji}\ge b^2$ and $q_{ij}\le R$, and off the
edges $q_{ij}=0$. Reversibility is the closed condition of Kolmogorov on cycle products,
and $\|S\|_2\le R$ is closed since $\pi$ is continuous. So $\mathcal P$ is compact, and
every connected face is irreducible for every $Q\in\mathcal C$.

\emph{Joint regularity.} The Perron root is simple, so $\lambda_F(v)$, the normalized
Perron vectors and $\alpha^*_F(v)$ are real-analytic in $(Q,v)$. The proof of
Proposition~\ref{prop:constant}(a), run with $Q$ as an extra variable, makes
$\theta(\alpha)$, $r_\alpha$, $l_\alpha$, $\Sigma_\alpha$, $J_F$ and $C_F$ real-analytic in
$(\alpha,Q)$, and $C_F$ is linear in $\mu$. So $K_F(v)$ is continuous and positive on
$\mathcal P$, which gives its bounds, and there is $\varepsilon_0>0$ with
$\alpha^*_{F,i}(v)\ge2k\varepsilon_0$ on $\mathcal P$. Put
$\mathcal K_0=\{\alpha\in\Delta_F:\alpha_i\ge\varepsilon_0\ \forall i\}$.

\emph{The constants.} They are the following, and each is continuous on $\mathcal P$ or on
$\mathcal P\times\mathcal K_0$, or comes from a compactness argument run there.
(i) The constant $C_\theta$ of Lemma~\ref{lem:upper} is built from the function $b$ of
Lemma~\ref{lem:torus}, which is continuous where the radius condition~\eqref{eq:A-radii}
holds strictly, as it does on $V_F$. So it is continuous along the Perron tilt.
(ii) In Theorem~\ref{thm:local} the Perron root in~\eqref{eq:A-radii} at $\theta'(\alpha)$ is
continuous and negative on $\mathcal P\times\mathcal K_0$, so one radius margin serves, and
Lemma~\ref{lem:laplace} applies on the compact space $\mathcal P\times\mathcal K_0$, since
the matrix $\mathcal J$ of Appendix~\ref{app:torus} is continuous and positive definite and
$\chi(\alpha,0)=(\mu\cdot r_\alpha)(l_\alpha\cdot\ones)/\alpha_1$ is continuous and positive.
(iii) In Step~3 of the proof of Lemma~\ref{lem:discrete} the lower bound for $H$ comes from
one word $c^0$ with at most $2k-1$ letters that walks along edges of $G$ through every state of
$F$. Its weight is at least $\mu_*(b^2/R)^{|c^0|-1}$, which bounds $H(T\alpha)$ below uniformly
for $T\ge1$ and $\alpha\in\mathcal K_0$. Step~4 uses only $q_{\max}\le R$ and
$n_{\min}\ge\varepsilon_0N$.
(iv) In Theorem~\ref{thm:facemax} the bounds for $\Sigma_\alpha^{-1}$, for the third
derivatives of $J_F$ and for two derivatives of $\log C_F$ are continuous on
$\mathcal P\times\mathcal K_0$, which gives one radius $\rho_0\le\varepsilon_0$ and uniform
$c_U$, $C_U$, $C_1$ and $C_3$. Away from the ball, $J_F(\beta)-v\cdot\beta-\lambda_F(v)$ is
lower semicontinuous and positive on a compact subset of $\mathcal P\times\Delta_F$, hence at
least $3\eta$, and finitely many vectors $\theta_u=A_Fu/u\in V_F$ serve, with constants
bounded by (i). For bounded $T$, $f_T$ is continuous and positive on
$[1,T_2]\times\mathcal P\times\mathcal K_0$.
\end{proof}

\begin{theorem}[finite-horizon recovery on the class]\label{thm:finite-horizon}
Run the design of Theorem~\ref{thm:reversible-reconstruction} at horizon $N$,
with $L=\log N$ and $T=\tau L$, and let $\mathcal W_N(\tau,v)$ be the set of
supports of the maximizers of $n\mapsto P_N(n)e^{Tv\cdot n/N}$. For an ordered
pair $(i,j)$ with $\{i,j\}\in E$ put $v=v^{ij}(\bar\delta)$,
\[
 \tau^-_N=\sup\Bigl\{\tau\in[\tfrac1{4R},U]:\mathcal W_N(\tau',v)=\{\{i\}\}\text{ for all }\tau'\in[\tfrac1{4R},\tau]\Bigr\},
\]
\[
 \tau^+_N=\inf\Bigl\{\tau\in[\tfrac1{4R},U]:\mathcal W_N(\tau',v)=\{\{i,j\}\}\text{ for all }\tau'\in[\tau,U]\Bigr\},
\]
let $t_{ij,N}$ be any number in $[\tau^-_N,\tau^+_N]$, and let
$K_{ij}=K_{\{i,j\}}(v)$. Put $\varepsilon_N=\log L/(2L)$. There are $N_0$ and $C$,
depending only on $G$, $R$, $b$ and $\mu_*$, such that for all $Q\in\mathcal C$,
$\mu\in\mathcal M$ and $N\ge N_0$ the following hold.
\begin{enumerate}[(a)]
\item (Singleton switches.) Let $i\ne1$ and $s=\tau_0L/N$. For every
$\delta\in[-3R,3R]$ other than
\[
 x_{i,N}=\frac1{\tau_0L}\Bigl[\log\frac{\mu_1}{\mu_i}+(N-1)\log\frac{1-sq_1}{1-sq_i}\Bigr],
\]
$\mathcal W_N(\tau_0,v^{i1}(\delta))$ is $\{\{1\}\}$ for $\delta<x_{i,N}$ and
$\{\{i\}\}$ for $\delta>x_{i,N}$. Moreover
$|x_{i,N}-x_i-2R\log(\mu_1/\mu_i)/L|\le CL/N$, so
$|x_{i,N}-x_i|\le2R\log(1/\mu_*)/L+CL/N$.
\item (Pair entries.) We have $\tau^+_N-\tau^-_N\le C/L^2$ and
\[
 \Bigl|t_{ij,N}-t_{ij}(1-\varepsilon_N)-\frac{c_{ij}}L\Bigr|\le C\,\frac{\log L}{L^2},\qquad
 c_{ij}=-t_{ij}\Bigl(\frac12\log t_{ij}+\log\frac{K_{ij}}{\mu_i}\Bigr),\qquad|c_{ij}|\le C .
\]
\item (Reconstruction.) Let $\widehat Q_N$ be the reconstruction of
Theorem~\ref{thm:reversible-reconstruction}(c) from the values $x_{i,N}$ and
$t_{ij,N}$, and $\widetilde Q_N$ the one from $x_{i,N}$ and
$t_{ij,N}/(1-\varepsilon_N)$. Both are reversible generators on $G$, and
\[
 \|\widehat Q_N-Q\|_2\le C\,\frac{\log L}L,\qquad\|\widetilde Q_N-Q\|_2\le\frac CL .
\]
\end{enumerate}
\end{theorem}

\begin{proof}
All constants below depend only on $G$, $R$, $b$ and $\mu_*$. On the ranges of
the design $T\in[L/(4R),UL]$, so for $N\ge N_1$ we have $TR\le N/2$,
$T\le N^{1/3}$ and $1/T+T^2/N\le5R/L$, and
Proposition~\ref{prop:uniform-facemax} applies.

\emph{Suppression.} Fix a pair and a point of the ranges, and let $m^*$ be the
least score among the irreducible faces inside the pair. The proof of
Lemma~\ref{lem:suppression} uses only $R$, $b$ and $\bar v$, so each connected
face $G'$ outside the pair has $\Phi_{\tau,v}(G')\ge m^*+\frac34$, and
Proposition~\ref{prop:uniform-facemax}(a) or the singleton atom gives
$M_{G'}\le CT^{|G'|-1}N^{-\Phi_{\tau,v}(G')}$. The best face inside the pair has
face maximum at least $c\,N^{-m^*}$, by the singleton atom with $\mu_i\ge\mu_*$
or by Proposition~\ref{prop:uniform-facemax}(b). Faces that are not connected
carry no mass. So for $N\ge N_0$ every maximizer has its support inside the
pair.

\emph{(a)} At $\tau_0$ the pair $\{i,1\}$, if it is an edge, has score at least
the better singleton score plus $\frac12$, by the proof of
Theorem~\ref{thm:reversible-reconstruction}(a), so as above it is smaller than
$\max(M_{\{1\}},M_{\{i\}})$ by a factor $CLN^{-1/2}$. The two singleton
atoms are exact, and
\[
 \log M_{\{i\}}-\log M_{\{1\}}=\log\frac{\mu_i}{\mu_1}+(N-1)\log\frac{1-sq_i}{1-sq_1}+\tau_0L\,\delta ,
\]
which is affine and increasing in $\delta$ with zero $x_{i,N}$. This gives the
claim on $\mathcal W_N$. Since $|\log(1-y)+y|\le y^2$ for $0\le y\le\frac12$,
the middle term is $-\tau_0L(q_i-q_1)+O(L/N)+O(R^2L^2/N)$, and dividing by
$\tau_0L=L/(2R)$ gives the expansion. For large $N$ it gives
$x_{i,N}\in(-3R,3R)$, since $|x_i|\le R$.

\emph{(b)} Let $\{i,j\}$ be an edge and $v=v^{ij}(\bar\delta)$. The exact atoms give
$\log M_{\{j\}}-\log M_{\{i\}}\le\log(1/\mu_*)-TR+1$ with $TR\ge L/4$, so $\{j\}$ loses to
$\{i\}$. By the proof of Theorem~\ref{thm:reversible-reconstruction}(b),
$\lambda_{\{i\}}(v)-\lambda_{\{i,j\}}(v)=1/t_{ij}$, and the exact atom gives
$\log M_{\{i\}}=\log\mu_i-T(q_i-\bar\delta)+O(1/L)$. So
Proposition~\ref{prop:uniform-facemax}(b) with $k=2$ gives, on the whole range,
\[
 D_N(T)=\log M_{\{i,j\}}-\log M_{\{i\}}=-L+\Xi(T)+\kappa+E_N(T),\qquad|E_N(T)|\le\frac{C_5}L,
\]
with $\Xi(T)=T/t_{ij}+\frac12\log T$ and $\kappa=\log(K_{ij}/\mu_i)$. Here
$|\kappa|\le C$, since $K_{ij}$ is bounded above and below and $\mu_i\ge\mu_*$.
The function $\Xi$ increases with $\Xi'\ge1/t_{ij}\ge2/U$. Let $T^\circ$ solve
$\Xi(T^\circ)=L-\kappa$. Then $D_N(T)<0$ for $T\le T^\circ-C_5U/L$ and
$D_N(T)>0$ for $T\ge T^\circ+C_5U/L$. As $t_{ij}\in[1/R,U/2)$, the window
$|T-T^\circ|<C_5U/L$ lies inside the range for large $N$. With the suppression
step, every maximizer has support $\{i\}$ below the window and $\{i,j\}$ above
it. Hence $T^\circ/L-C_5U/L^2\le\tau^-_N\le\tau^+_N\le T^\circ/L+C_5U/L^2$,
which is the first claim, and $t_{ij,N}=T^\circ/L+O(1/L^2)$. To solve for
$T^\circ$ put $\tau^\circ=T^\circ/L$. Then
$\tau^\circ=t_{ij}\bigl(1-(\frac12\log T^\circ+\kappa)/L\bigr)$, so
$\tau^\circ/t_{ij}=1+O(\log L/L)$ and
$\log T^\circ=\log L+\log t_{ij}+O(\log L/L)$. Putting this back gives
$\tau^\circ=t_{ij}(1-\varepsilon_N)+c_{ij}/L+O(\log L/L^2)$, uniformly, since
$t_{ij}\le U/2$ and $|\kappa|\le C$.

\emph{(c)} By (a) and (b) the dimensionless coordinates $x_{i,N}/R$ and
$Rt_{ij,N}$ differ from $x_i/R$ and $Rt_{ij}$ by at most $C\log L/L$, and after
the correction by at most $C/L$, since
$t_{ij,N}/(1-\varepsilon_N)-t_{ij}=(c_{ij}/L+O(\log L/L^2))/(1-\varepsilon_N)$.
Theorem~\ref{thm:stability} gives its admissible error size $\varepsilon_0$ and
its constant as functions of $d$ and $b/R$ alone. So for $N\ge N_0$ these errors
are below $\varepsilon_0$, and that theorem gives both bounds and the generator
property.
\end{proof}

\begin{remark}[a start that defeats every horizon]\label{rem:finite-start}
The bound $\min_i\mu_i\ge\mu_*$ cannot be dropped, since the start enters
through the singleton switches as $2R\log(\mu_1/\mu_i)/L$. Let $d=2$, $q_{12}=q_{21}=1$, $\tau=\frac14$ and $v=(0,1)$. The scores are
$\Phi(\{2\})=0$, $\Phi(\{1\})=\frac14$ and
$\Phi(\{1,2\})=1+\frac14\cdot\frac{1-\sqrt5}2>\frac34$, so for a fixed start
$\{2\}$ wins for large $N$. Both singleton atoms carry the factor
$(1-T/N)^{N-1}$, so $M_{\{1\}}/M_{\{2\}}=(\mu_1/\mu_2)e^{-T}=(\mu_1/\mu_2)N^{-1/4}$
exactly. For the start $\mu_2=1/(1+2N^{1/4})$ this ratio is $2$ at every $N$,
while Lemma~\ref{lem:upper} with the Perron tilt, whose constant is bounded
uniformly in the start, gives $M_{\{1,2\}}\le CLN^{-1/2}M_{\{1\}}$. So $\{1\}$
wins at every large $N$, against the limiting prediction. The switch between
$\{1\}$ and $\{2\}$ in the variable $v_2$ sits at
$\log(\mu_1/\mu_2)/(\tau L)=1+4\log2/L$ instead of at $q_2-q_1=0$. On
$\mathcal M$ the same shift is at most $\log(1/\mu_*)/(\tau L)$, uniformly.
\end{remark}

\begin{remark}[the size of the correction]\label{rem:finite-size}
On the two-state class with $R=3$ and $b=1$, the chain $q_1=q_2=1$ with the
uniform start has $t_{12}=6.16$ and $c_{12}=-16.5$, and at $N=10^7$ the exact law
gives $t_{12,N}=4.65$, against $5.63$ for $t_{12}(1-\varepsilon_N)$ and $4.60$ for
the prediction of (b), so the term $c_{ij}/L$, which no correction fixed in
advance removes, dominates $t_{ij}\varepsilon_N$ at accessible $N$.
\end{remark}
